\fig{fv:s2:minimum}{允许射线上的最近点：垂足在射线上与垂足在射线外}{\begin{tikzpicture}

\begin{scope}
\coordinate(O)at(0,0);\coordinate(K)at(0,2.4);\coordinate(H)at(-1.2,1.2);
\draw[axis](O)--(-2.1,2.1)node[left]{$\lambda\bm e,\ \lambda\ge0$};\draw[flow](O)--(K)node[right]{$-\bm G$};
\draw[force](H)--(K)node[midway,above left]{$\bm F_{\min}$};\fill(H)circle(2pt)node[left]{$H$};\node[below]at(O){$O$};\node[above]at(K){$K$};
\draw(-1.06,1.06)--(-.92,1.2)--(-1.06,1.34);\node at(-.4,-.8){$\lambda_*=-\bm G\!\cdot\!\bm e>0$};
\end{scope}
\begin{scope}[xshift=5cm]
\draw[axis](0,0)--(-1.7,1.7)node[left]{$\lambda\bm e$};\draw[guide](0,0)--(1.4,-1.4);\draw[force](0,0)--(1,-1)node[right]{$K=-\bm G$};
\fill(0,0)circle(2pt)node[above right]{$O=H_{\rm feasible}$};\node[align=center]at(.2,2.2){垂足在反向延长线上};\node at(.2,-1.8){可行最小值：$\lambda=0,\ F=G$};
\end{scope}
\end{tikzpicture}}
图\ref{fv:s2:minimum}先固定向量 $\overrightarrow{OK}=-\bm G$。反力端点只能在 $\lambda\bm e$ 的允许射线上移动，外力是从该端点指向 $K$ 的向量。左图最近点是垂足 $H$，所以最小外力垂直反力；右图虽然直线延长线上有更近的点，却对应负张力或负支持力，允许射线上的最近点只能是端点 $O$。作图时保留射线的箭头，便能同时看到极值条件和可实现条件。
